% LaTeX companion of hw09-quotient-rings.typ. Both formats are editable.
\chapter{Homework 9}

\section{1. Prime-order groups}

(a) Prove Fermat's Little Theorem: if \(p\) is prime and \(\left. p! \middle| \ a \right.\), then \(a^{p - 1} \equiv 1\ \text{mod}\ p\).

(b) If \(G\) is a group of prime order \(p\), then \(G\) is cyclic.

(c) A nontrivial group \(G\) has no nontrivial proper subgroups if and only if \(G\) is finite and of order \(p\) where \(p\) is prime.

\subsection{(a)}

\textbf{Claim 1.} If \(p\) is prime, then \(a^{p} \equiv a\ \text{mod}\ p\).

\textbf{Proof.} Take arbitrary prime \(p\). We prove it by induction on \(a\).

\textbf{Basic step.} \(1^{p} = 1\), so \(1^{p} \equiv 1\ \text{mod}\ p\).

\textbf{Inductive step.} Assume \(a^{p} \equiv a\ \text{mod}\ p\). We show that \(\left( {a + 1} \right)^{p} \equiv \left( {a + 1} \right)\ \text{mod}\ p\). By the binomial theorem,

\[\begin{matrix}
\left( {a + 1} \right)^{p} & {= \sum\limits_{k = 0}^{p}\binom{p}{k}a^{k}} \\
 & {= \sum\limits_{k = 0}^{p}\frac{p!}{k!\left( {p - k} \right)!}a^{k}} \\
 & {= \sum\limits_{k = 1}^{p - 1}\frac{p!}{k!\left( {p - k} \right)!}a^{k} + 1 + a^{p}.}
\end{matrix}\]

For every \(1 \leq k \leq p - 1\), \(p - k \leq p - 1\) and \(p - k \geq 1\), so \(p\) divides every term with denominator \(k!\left( {p - k} \right)!\): since \(p\) is prime, \(\left. p! \middle| \ k!\left( {p - k} \right)! \right.\), otherwise \(p\) must divide one of the factors in that product, which contradicts the bounds. Therefore

\[\frac{\left( {p - 1} \right)!}{k!\left( {p - k} \right)!}\]

is still an integer, and

\[\left( {a + 1} \right)^{p} = p\left( {\sum\limits_{k = 1}^{p - 1}\frac{\left( {p - 1} \right)!}{k!\left( {p - k} \right)!}a^{k}} \right) + 1 + a^{p}.\]

Thus \(p\) divides \(\left( {a + 1} \right)^{p} - \left( {a^{p} + 1} \right)\), and therefore

\[\left( {a + 1} \right)^{p} \equiv a^{p} + 1 \equiv a + 1\ \text{mod}\ p.\]

This proves Claim 1.

\textbf{Claim 2.} Following Claim 1, if \(\left. p! \middle| \ a \right.\), then \(a^{p - 1} \equiv 1\ \text{mod}\ p\).

\textbf{Proof.} Let \(p\) be an arbitrary prime and take arbitrary \(a \in \mathbb{Z}^{+}\) with \(\left. p! \middle| \ a \right.\). By Claim 1,

\[a^{p} \equiv a\ \text{mod}\ p,\]

so \(\left. p\  \middle| \ a\left( {a^{p - 1} - 1} \right) \right.\). Since \(p\) is prime, either \(\left. p\  \middle| \ a \right.\) or \(\left. p\  \middle| \ a^{p - 1} - 1 \right.\). Since \(\left. p! \middle| \ a \right.\), we get

\[a^{p - 1} \equiv 1\ \text{mod}\ p.\]

Combining Claims 1 and 2 proves Fermat's Little Theorem.

\subsection{(b)}

\textbf{Proof.} Assume \(|G|\) is prime, so \(|G| \geq 2\) and there is a non-identity element in \(G\). Select arbitrary non-identity \(a \in G\). Then \(\left| \left\langle a \right\rangle \right| \geq 2\), since \(a \land a^{2} \in \left\langle a \right\rangle\) (with \(a \neq e\), otherwise \(aa^{- 1} = a^{- 1}a\) would give \(a = e\)). By Lagrange's Theorem,

\[|G| = \left| \left\langle a \right\rangle \right| \cdot \text{index of}\ \left\langle a \right\rangle\ \text{in}\ G.\]

Since \(|G|\) is prime and \(\left| \left\langle a \right\rangle \right| \geq 2\), we have \(\left| \left\langle a \right\rangle \right| = |G|\), which means \(\left\langle a \right\rangle = G\). Hence \(G\) is cyclic.

\subsection{(c)}

First we prove the backward direction. Assume \(|G|\) is finite and prime. Then for every subgroup \(K \leq G\), Lagrange's Theorem gives \(\left. |K|\  \middle| \ |G| \right.\). Since \(|G|\) is prime, \(|K| = 1\) or \(|G|\), so \(K\) is either trivial or \(G\) itself. Therefore \(G\) has no nontrivial subgroups.

For the forward direction, assume \(G\) has no nontrivial proper subgroup. Case 1: \(G\) has finite composite order. Then \(|G| = mn\) for some prime \(m\) and \(n \geq 2\). By Theorem 8.6, for every \(x \in G\), \(x^{|G|} = e\). Pick \(x \neq e\). If \(x^{m} = e\), then \(\left\langle x \right\rangle\) has order at most \(m < |G|\) and is nontrivial, a contradiction. If \(x^{m} \neq e\), then \(\left\langle x^{m} \right\rangle\) has order at most \(n < |G|\) and is nontrivial, another contradiction.

Case 2: \(G\) has infinite order. Select arbitrary non-identity \(g \in G\) and consider \(\left\langle g \right\rangle\). If \(g \in \left\langle g^{2} \right\rangle\), then \(g = \left( g^{2} \right)^{n} = g^{2n}\) for some integer \(n\), so \(g^{2n - 1} = e\) and \(\left| \left\langle g \right\rangle \right| \leq 2n - 1\); this is a nontrivial proper subgroup of \(G\), a contradiction. If \(g! \in \left\langle g^{2} \right\rangle\), then \(\left\langle g^{2} \right\rangle\) is itself a nontrivial proper subgroup of \(G\), again a contradiction. Thus the group cannot be infinite. It must have prime order.

\section{2. Left and right cosets}

For each of the following parts, \(K\) is a subgroup of the group \(G\). Write down every element of every distinct right coset and every distinct left coset.

(a) \(K = \left\{ {,r_{0}\text{,}r_{90}\text{,}r_{180}\text{,}r_{270},} \right\}\) and \(G = D_{4}\), with reflections \(s_{v},s_{h},s_{\text{NW}},s_{\text{NE}}\).

(b) \(K = \left\{ {,e\text{,}\left( {1,2} \right),} \right\}\) and \(G = S_{3}\).

(c) \(K = \left\langle \begin{pmatrix}
0 & 1 \\
1 & 0
\end{pmatrix} \right\rangle\) and \(G = \operatorname{GL}_{2}\left( \mathbb{Z}_{2} \right)\), whose elements are

\[\left\{ {,I = \begin{pmatrix}
1 & 0 \\
0 & 1
\end{pmatrix}\text{,}a = \begin{pmatrix}
0 & 1 \\
1 & 0
\end{pmatrix}\text{,}b = \begin{pmatrix}
1 & 1 \\
1 & 0
\end{pmatrix}\text{,}c = \begin{pmatrix}
1 & 1 \\
0 & 1
\end{pmatrix}\text{,}d = \begin{pmatrix}
1 & 0 \\
1 & 1
\end{pmatrix}\text{,}f = \begin{pmatrix}
0 & 1 \\
1 & 1
\end{pmatrix},} \right\}.\]

(d) \(K = \left\langle 5 \right\rangle\) and \(G = \mathbb{Z}_{12}^{\times}\).

\subsection{(a)}

There are two left/right cosets.

\textbf{Left cosets:}

\begin{enumerate}
\item
  \(r_{0}K = r_{90}K = r_{180}K = r_{270}K = K\).
\item
  \(s_{v}K = \left\{ {,s_{v}\text{,}s_{h}\text{,}s_{\text{NW}}\text{,}s_{\text{NE}},} \right\} = s_{h}K = s_{\text{NW}}K = s_{\text{NE}}K\).
\end{enumerate}

\textbf{Right cosets:}

\begin{enumerate}
\item
  \(Kr_{0} = Kr_{90} = Kr_{180} = Kr_{270} = K\).
\item
  \(Ks_{v} = Ks_{h} = Ks_{\text{NW}} = Ks_{\text{NE}} = \left\{ {,s_{v}\text{,}s_{h}\text{,}s_{\text{NW}}\text{,}s_{\text{NE}},} \right\}\).
\end{enumerate}

\subsection{(b)}

\[G = S_{3} = \left\{ {,e\text{,}\left( {1,2} \right)\text{,}\left( {1,3} \right)\text{,}\left( {2,3} \right)\text{,}\left( {1,2,3} \right)\text{,}\left( {1,3,2} \right),} \right\}.\]

There are three left/right cosets.

\textbf{Left cosets:}

\begin{enumerate}
\tightlist
\item
  \(\left( {1,3} \right)K = \left\{ {,\left( {1,3} \right)\text{,}\left( {1,2,3} \right),} \right\}\).
\item
  \(\left( {2,3} \right)K = \left\{ {,\left( {2,3} \right)\text{,}\left( {1,3,2} \right),} \right\}\).
\item
  \(K\) itself.
\end{enumerate}

\textbf{Right cosets:}

\begin{enumerate}
\tightlist
\item
  \(K\left( {1,3} \right) = \left\{ {,\left( {1,3} \right)\text{,}\left( {1,3,2} \right),} \right\}\).
\item
  \(K\left( {2,3} \right) = \left\{ {,\left( {2,3} \right)\text{,}\left( {1,2,3} \right),} \right\}\).
\item
  \(K\) itself.
\end{enumerate}

\subsection{(c)}

\(K = \left\{ {,a\text{,}I,} \right\}\). There are three left/right cosets.

\textbf{Left cosets:}

\begin{enumerate}
\tightlist
\item
  \(aK = IK = K = \left\{ {,a\text{,}I,} \right\}\).
\item
  \(fK = \left\{ {,f\text{,}d,} \right\} = dK\).
\item
  \(bK = \left\{ {,c\text{,}b,} \right\} = cK\).
\end{enumerate}

\textbf{Right cosets:}

\begin{enumerate}
\tightlist
\item
  \(Ka = KI = K = \left\{ {,a\text{,}I,} \right\}\).
\item
  \(Kf = \left\{ {,c\text{,}f,} \right\} = Kc\).
\item
  \(Kd = \left\{ {,b\text{,}d,} \right\} = Kb\).
\end{enumerate}

\subsection{(d)}

\[G = \left\{ {,1\text{,}5\text{,}7\text{,}11,} \right\},\quad K = \left\langle 5 \right\rangle = \left\{ {,5\text{,}1,} \right\}.\]

There are two left/right cosets.

\textbf{Left cosets:} \(K = \left\{ {,5\text{,}1,} \right\}\) and \(7K = \left\{ {,11\text{,}7,} \right\}\).

\textbf{Right cosets:} \(K = \left\{ {,5\text{,}1,} \right\}\) and \(K7 = \left\{ {,11\text{,}7,} \right\}\).

\section{3. Conjugacy classes}

Any group \(G\) acts on itself by conjugation: \(g \cdot h = ghg^{- 1}\). The orbits of this action are called \emph{conjugacy classes}.

\begin{enumerate}
\tightlist
\item
  Show \(h \in Z(G)\) if and only if \(h\) is a fixed point of the conjugation action.
\item
  Show a subgroup \(H\) of \(G\) is normal if and only if it is a disjoint union of conjugacy classes.
\item
  Describe the partition of \(S_{5}\) into its conjugacy classes.
\item
  Show that the only nontrivial normal subgroup of \(S_{5}\) is \(A_{5}\).
\end{enumerate}

\subsection{1.}

\textbf{Forward direction.} Assume \(h\) is a fixed point of the conjugation action. Then for every \(g \in G\), \(ghg^{- 1} = h\). Multiply by \(g\) on both sides to get \(gh = hg\), so \(h \in Z(G)\).

\textbf{Backward direction.} Assume \(h \in Z(G)\). Then for every \(g \in G\), \(gh = hg\), so \(ghg^{- 1} = hgg^{- 1} = h\). Hence \(h\) is a fixed point of the conjugation action. Thus \(h\) is a fixed point of the conjugation action if and only if \(h \in Z(G)\).

\subsection{2.}

\textbf{Forward direction.} Assume \(H\) is a disjoint union of conjugacy classes, i.e.

\[H = \cup_{j = 1}^{n}\left\{ ,gh_{j}g^{- 1}\  \middle| \ g \in G, \right\}\]

for some \(h_{1},h_{2},\ldots,h_{n} \in G\). Select arbitrary \(g \in G\) and fix it. Take arbitrary \(x \in H\); then \(x \in O\left( h_{i} \right)\) for some \(h_{i}\), so \(gxg^{- 1} = g \cdot x \in O\left( h_{i} \right) \subset H\). Thus \(g^{- 1}xg \in H\) and, for every \(g \in G\), \(gHg^{- 1} \subset H\). By Theorem 8.11, \(H\) is a normal subgroup of \(G\).

\textbf{Backward direction.} Assume \(H\) is normal. Take arbitrary \(g \in G\). By Theorem 8.11, \(gHg^{- 1} \subset H\), so for every \(h \in H\), \(ghg^{- 1} \in H\). Since \(g\) is arbitrary, \(O(h) \subset H\) for every \(h \in H\). Therefore \(H\) is a union of conjugacy classes. Since orbits are either disjoint or identical, it is a disjoint union of conjugacy classes.

\subsection{3.}

The conjugacy classes in \(S_{5}\) are:

\[\begin{matrix}
{O(e)} & {= \left\{ {,e,} \right\}\ } & & {\text{order}\ 1} \\
{O\left( \left( {1,2} \right) \right)} & {= \left\{ {,\text{all 2-cycles},} \right\}\ } & & {\text{order}\ \binom{5}{2} = 10} \\
{O\left( \left( {1,2,3} \right) \right)} & {= \left\{ {,\text{all 3-cycles},} \right\}\ } & & {\text{order}\ \binom{5}{3} \times 2 = 20} \\
{O\left( \left( {1,2,3,4} \right) \right)} & {= \left\{ {,\text{all 4-cycles},} \right\}\ } & & {\text{order}\ \binom{5}{4} \times 3 \neq 30} \\
{O\left( \left( {1,2,3,4,5} \right) \right)} & {= \left\{ {,\text{all 5-cycles},} \right\}\ } & & {\text{order}\ 4 \neq 24} \\
{O\left( {\left( {1,2} \right)\left( {3,4} \right)} \right)} & {= \left\{ {,\text{all two disjoint transpositions},} \right\}\ } & & {\text{order}\ \left( \frac{1}{2} \right)\binom{5}{2}\binom{3}{2} = 15} \\
{O\left( {\left( {1,2} \right)\left( {3,4,5} \right)} \right)} & {= \left\{ {,\text{all 2+3-disjoint cycles},} \right\}\ } & & {\text{order}\ \binom{5}{3} \times 2 = 20.}
\end{matrix}\]

These union to \(S_{5}\), with orders summing to \(120\).

\subsection{4.}

Let \(K\) be a nontrivial normal subgroup of \(S_{5}\). First \(e \in K\) by the definition of subgroup. By part 2, \(K\) is a disjoint union of conjugacy classes, so \(\left\{ {,e,} \right\}\) is one of the conjugacy classes that form \(K\). By Lagrange's Theorem, \(\left. |K|\  \middle| \ \left| S_{5} \right| = 120 \right.\).

Since \(K \neq \left\{ {,e,} \right\}\), more conjugacy classes must be in the disjoint union. Since \(\left| \left\{ {,e,} \right\} \right| = 1\) and the class of all 5-cycles has order \(24\), the all-5-cycles class must be one of the conjugacy classes; otherwise \(|K|\) cannot divide \(120\). Now \(|K| \geq 25\). Thus \(|K|\) can only be \(30,40,\) or \(60\) to divide \(120\), and all two-disjoint transpositions of order \(15\) must be one of the classes.

There are three possibilities:

\begin{enumerate}
\tightlist
\item
  \(K = \left\{ {,e,} \right\} \cup \left\{ {,\text{all 5-cycles},} \right\} \cup \left\{ {,\text{all two disjoint transpositions},} \right\}\).
\item
  The preceding union together with \(\left\{ {,\text{all 3-cycles},} \right\}\).
\item
  The preceding union together with \(\left\{ {,\text{all 2+3-disjoint cycles},} \right\}\).
\end{enumerate}

A normal subgroup must be closed under operation and inverse. For case 1,

\[\left( {3,4} \right)\left( {1,2} \right)\left( {1,2,3,4,5} \right) = \left( {2,4,5} \right)! \in K,\]

so it is not a subgroup. For case 3,

\[\left( {3,4} \right)\left( {1,2,3} \right)\left( {1,2,3,4,5} \right) = \left( {1,4,5,2} \right)! \in K,\]

so it is not a subgroup. Therefore only case 2 can be a subgroup. Its elements are even, so it is \(A_{5}\). Thus \(K\) is the only nontrivial normal subgroup of \(S_{5}\).

\section{\texorpdfstring{4. A group whose order is divisible by \(p\)}{4. A group whose order is divisible by p}}

Let \(p\) be a prime, and \(G\) a finite group with \(\left. p\  \middle| \ |G| \right.\). Consider

\[X = \left\{ ,\left( {g_{1},\ldots,g_{p}} \right) \in G \times \ldots \times G\  \middle| \ g_{1}g_{2}\cdots g_{p} = e, \right\},\]

where there are \(p\) copies of \(G\). The group \(\mathbb{Z}_{p}\) acts on \(X\) by rotating elements:

\[i_{p} \cdot \left( {g_{1},\ldots,g_{p}} \right) = \left( {g_{1 + i},\ldots,g_{p},g_{1},\ldots,g_{i}} \right).\]

\begin{enumerate}
\tightlist
\item
  Show \(X\) has \(|G|^{p - 1}\) elements, so \(\left. p\  \middle| \ |X| \right.\).
\item
  Show the orbits of the action either have \(1\) or \(p\) elements, and the orbits of order \(1\) are either \(\left( {e,e,\ldots,e} \right)\) or of the form \(\left( {g,g,\ldots,g} \right)\) with \(|g| = p\).
\item
  Show that \(G\) contains an element of order \(p\).
\end{enumerate}

\subsection{1.}

For any choice of \(\left( {g_{1},g_{2},\ldots,g_{p - 1}} \right)\), by existence and uniqueness of inverses there is a fixed \(g_{p}\) such that \(g_{1}g_{2}\cdots g_{p} = e\). Therefore there are \(|G|^{p - 1}\) choices of \(g_{1},\ldots,g_{p - 1}\) and

\[|X| = |G|^{p - 1} = |G||G|^{p - 2},\]

so \(\left. p\  \middle| \ |X| \right.\).

\subsection{2.}

For \(x \in X\), the Orbit-Stabilizer Theorem gives

\[\left| {\operatorname{Orbit}(x)} \right| \cdot \left| {\operatorname{Stab}(x)} \right| = \left| \mathbb{Z}_{p} \right| = p.\]

Consider \(x = \left( {g,g,g,\ldots,g} \right)\) for some \(g \in G\). For every \(y \in \mathbb{Z}_{p}\), \(y \cdot x = x\), so \(\operatorname{Stab}(x) = \mathbb{Z}_{p}\); therefore \(\left| {\operatorname{Stab}(x)} \right| = p\) and \(\left| {\operatorname{Orbit}(x)} \right| = 1\). In this case either \(g = e\) or \(|g| = p\), because \(\left( {g,g,\ldots,g} \right) \in X\) means \(g^{p} = e\). Thus either \(g = e\) or \(|g| = p\); \(|g|\) cannot be less than \(p\), since otherwise \(g^{p} = g^{{|g|}a} = e\) for some \(a \in \mathbb{Z}\) would contradict that \(p\) is prime.

Otherwise, in \(x = \left( {g_{1},g_{2},\ldots,g_{p}} \right)\) at least some \(g_{i},g_{j}\) are different. Only when \(y = 0_{p}\) does \(y \cdot x = x\), so \(\operatorname{Stab}(x) = \left\{ {,0_{p},} \right\}\) and \(\left| {\operatorname{Orbit}(x)} \right| = p\). Therefore the orbits of the action of \(\mathbb{Z}_{p}\) on \(X\) either have \(1\) or \(p\) elements, and the orbits of order \(1\) are either \(\left( {e,e,\ldots,e} \right)\) or \(\left( {g,g,\ldots,g} \right)\) with \(|g| = p\).

\subsection{3.}

Every element of \(X\) belongs to exactly one orbit. From part 2,

\[|X| = mp = np\]

for some \(m \in \mathbb{Z}\) after reducing the \(p\)-element orbits. Since \(\left. p\  \middle| \ |X| \right.\), there must be an element \(\left( {g,g,\ldots,g} \right) \in X\) distinct from \(\left( {e,e,\ldots,e} \right)\) such that \(|g| = p\). Therefore there must exist \(g \in G\) such that \(|g| = p\).

\section{Source notes}

The handwritten proof for Problem 1(a) states the induction as ``on \(a\)'' after introducing a prime \(p\); the typeset version retains that stated induction. In Problem 4(3), the source uses the notation \(|X| = mp - np\) while reducing orbit counts; it is transcribed as written.
